\section{Alternative Methods and the Proof of Fubini in 3D}

\begin{nice-box}[Cartesian Slicing vs. 3D Box Splitting]
The left board contains a diagram of a 2D cross-section of a sphere sliced at height $x_3$. The right board sets up the formal 3D box definitions for the proof of Fubini's theorem.
\end{nice-box}

\begin{math_stroke}[Cartesian Slicing and Box Definition]
\textbf{Left Board: Cartesian Slicing of the Sphere} \\
Using the Pythagorean theorem, the radius of the slice is written as:
\[
r^2 + x_3^2 = 1 \implies r = \sqrt{1 - x_3^2}
\]
The Cavalieri volume integral is established:
\[
V = \int_{-1}^{1} \pi (1 - x_3^2) \, dx_3
\]

\textbf{Right Board: 3D Box Definition for Fubini's Proof}
\[
K = [a_1, b_1] \times [a_2, b_2] \times [a_3, b_3] \subset \mathbb{R}^3
\]
The splitting strategy ($k=1$):
\[
\mathbb{R}^3 \equiv \underbrace{\mathbb{R}^1}_{x_1} \times \underbrace{\mathbb{R}^2}_{(x_2, x_3)}
\]
\end{math_stroke}
\begin{spoken_clean}[15:00 - 18:38]
So, is it clear how to apply the theorem in practice? Is it useful? Yes, absolutely. Now, we will move on to prove the corollary we just used. We understand what this corollary says, and the proof itself is not particularly difficult, although the formal statement can be a bit intimidating to read at first. Because this is the kind of exercise that everyone needs to know how to do, we must go through it. Are there any questions before we begin the proof?
\end{spoken_clean}

\begin{student_question}
A student raises their hand and asks: "Would it not be easier to calculate the volume of the unit sphere just using the slicing lemma directly, without doing a change of variables? Just slicing the sphere into circles?"
\end{student_question}

\begin{spoken_clean}[18:38 - 21:05]
You are asking if we could calculate the volume of the unit sphere without doing a change of variables, simply by slicing it into circles. It is true, you can certainly slice the sphere like that. You can do it, but it is not necessarily easier! Let's look at what that entails. You can slice it horizontally, along the $x_3$-axis, which we can call the height. 

The Cavalieri principle says that you can integrate from $-1$ to $1$, and in each slice, you compute the area of the resulting disk. By using Pythagoras, if the hypotenuse is $1$, and the vertical distance is $x_3$, the radius squared of this disk is exactly $1 - x_3^2$. So, you can also evaluate this integral to find the volume. However, notice that you must remember how to compute the primitive of this specific polynomial function. You always need to do some extra work here. In contrast, with the polar coordinates we just used, the primitives were completely trivial.
\end{spoken_clean}

\begin{math_stroke}[Direct Cartesian Slicing]
Setting up the integral based on the student's suggestion:
\[
\text{Volume} = \underbrace{\int_{-1}^{1}}_{\text{``integrate from } -1 \text{ to } 1 \text{''}} \underbrace{\pi \big(1 - x_3^2\big)}_{\text{``the radius squared of this disk is exactly } 1 - x_3^2 \text{''}} \, \underbrace{dx_3}_{\text{``along the } x_3 \text{-axis''}}
\]
\end{math_stroke}

\begin{didactic_insight}[The Hidden Cost of Cartesian Integration]
Acknowledging geometric intuition highlights a vital pragmatic truth: Cartesian slicing often leads to integrating roots and complex polynomials. Polar coordinates, while requiring the initial setup of a Jacobian determinant, radically simplify the actual calculus, reducing the integrands to basic trigonometric functions and constants.
\end{didactic_insight}

\begin{spoken_clean}[21:05 - 24:15]
Okay, let us move on to the proof of the corollary. Because this is Analysis II, there is often a big leap from Analysis I—where you can visualize everything very well—to abstract $n$-dimensional spaces, where some people struggle. That is completely normal! 

For concreteness, I will do this proof in three dimensions. The proof in your notes is in $n$ dimensions, and it is exactly the same, but with a general $n$. This is an exercise I am telling you to do all the time: whenever a proof is too abstract in $n$ dimensions, just write it down in 2 or 3 dimensions, and you will see exactly what is going on. Simplify. How do we proceed in 3D? We will use Fubini's theorem with $k = 1$. We have our box, $K$, and we split the space such that the first variable is in one dimension, and the other two are in a two-dimensional box.
\end{spoken_clean}

\begin{math_stroke}[First Dimensional Split ($k=1$)]
Applying the general Fubini theorem to the 3D box $K$:
\[
\int_K f = \underbrace{\int_{a_1}^{b_1}}_{\text{``first variable is in one dimension''}} \underbrace{\Phi_1(x_1)}_{\text{``marginal function''}} \, dx_1
\]
\end{math_stroke}

\begin{spoken_clean}[24:15 - 26:10]
The integral of our function over the box is equal to the integral over $[a_1, b_1]$ of the marginal function, which we will call $\Phi_1(x_1)$. Now, what exactly is this marginal? It is the integral over the remaining two-dimensional box, $[a_2, b_2] \times [a_3, b_3]$, of our function with respect to $x_2$ and $x_3$. 

Because $f$ is continuous over the entire box, if we fix $x_1$, the restriction of $f$ to this 2D plane is also a continuous function. Since it is a continuous function over a closed, two-dimensional box, it is automatically integrable. We do not need to worry about upper and lower integrals here; the integral is perfectly well-defined.
\end{spoken_clean}

\begin{math_stroke}[Expanding the Marginal $\Phi_1$]
\[
\Phi_1(x_1) = \underbrace{\int_{[a_2, b_2] \times [a_3, b_3]}}_{\text{``remaining two-dimensional box''}} \underbrace{f(x_1, x_2, x_3)}_{\text{``if we fix } x_1 \text{, the restriction is continuous''}} \, \underbrace{d(x_2, x_3)}_{\text{``with respect to } x_2 \text{ and } x_3 \text{''}}
\]
\end{math_stroke}

\begin{spoken_clean}[26:10 - 28:00]
Now, we need to compute this inner, two-dimensional integral. But we can simply apply the exact same idea again! To compute this integral, we use Fubini's theorem once more with $k = 1$, but this time, in dimension 2. We reduce the problem by one more dimension. 

We integrate over $[a_2, b_2]$ with respect to $x_2$, and the new marginal inside is the integral over $[a_3, b_3]$ of $f$ with respect to $x_3$. Here, both $x_1$ and $x_2$ are fixed. We integrate with respect to $x_3$, we get a number, and then we move $x_2$. 
\end{spoken_clean}

\begin{math_stroke}[Second Dimensional Split (Applying Fubini to the 2D slice)]
\[
\int_{[a_2, b_2] \times [a_3, b_3]} f \, d(x_2, x_3) = \underbrace{\int_{a_2}^{b_2}}_{\text{``integrate over } [a_2, b_2] \text{''}} \left( \underbrace{\int_{a_3}^{b_3} f(x_1, x_2, x_3) \, dx_3}_{\text{``both } x_1 \text{ and } x_2 \text{ are fixed''}} \right) \underbrace{dx_2}_{\text{``then we move } x_2 \text{''}}
\]
\end{math_stroke}

\begin{spoken_clean}[28:00 - 30:00]
If you substitute this all back into our first equation, you see that it gives you exactly the triple nested formula I wrote before: an integral from $a_1$ to $b_1$, $a_2$ to $b_2$, and $a_3$ to $b_3$. End of the proof! 

If you look at the notes, it is the exact same proof in $n$ dimensions. It simply introduces notation for these successive domains that become smaller and smaller as we integrate out variables one by one. There are two things in every theorem: the theoretical proof, and the practical application. To pass the exam, you need to at least be able to apply the theorem without hesitation. If you have some doubts about the abstract proof, that understanding will come with time, but understanding how to apply it is absolutely crucial.
\end{spoken_clean}

\begin{nice-box}
\begin{theorem}[Concluded 3D Fubini Iteration]
Substituting the nested marginals back into the original split yields the final iterated form:
\[
\int_K f = \underbrace{\int_{a_1}^{b_1} \int_{a_2}^{b_2} \int_{a_3}^{b_3} f(x_1, x_2, x_3) \, dx_3 \, dx_2 \, dx_1}_{\text{``exactly the triple nested formula I wrote before''}}
\]
This recursive logic extends to $n$ dimensions.
\end{theorem}
\end{nice-box}